\capitulo{2}{Objetivos}

\section{Objetivo general}
Integrar conocimientos adquiridos durante el grado en un proyecto completo
 que abarca desde la preparación de los datos hasta el desarrollo de una 
 aplicación funcional. La cual, consiste en una herramienta, capaz de detectar posibles interacciones farmacológicas entre principios activos, 
mediante el análisis de las enzimas que las metabolizan, además de mostrar recomendaciones de fármacos con el mismo efecto 
terapéutico pero sin riesgo de interacción.

\section{Objetivos específicos}
\begin{enumerate}
    \item Analizar las principales bases de datos farmacológicas disponibles para la obtención de información sobre metabolismo, interacciones y efectos adversos de los medicamentos.
    \item Diseñar un proceso de extracción, transformación y normalización de datos heterogenéos que permita integrar la información procedentes de diferentes bases de datos farmacológicas en una base de datos unificada.
    \item Resolver ambigüedades en aquellos principios activos para los que se detecta mas de una vía metabólica principal.
    \item Desarrollar un algoritmo capaz de identificar posibles interacciones entre dos principios activos a partir de las enzimas implicadas en su metabolismo.
    \item Implementar un sistema de clasificación del nivel de riesgo de interacción basado en la coincidencia de enzimas metabólicas principales y secundarias.
    \item Incorporar información complementaria sobre efectos adversos y alternativas terapéuticas mediante la utilización de códigos ATC.
    \item Desarrollar e implementar una interfaz intuitiva que facilite la consulta y visualización de los resultados obtenidos.
    \item Validar el correcto funcionamiento mediante diferentes casos de prueba.
\end{enumerate}

\section{Objetivos de aprendizaje}
\begin{enumerate}
    \item Implementación para el desarrollo de aplicaciones de análisis de datos mediante la profundización del manejo del lenguaje de programación Python.
    \item Adquirir experiencia en el desarrollo de aplicaciones web utilizando el framework Dash.
    \item Profundizar en el manejo de técnicas de procesamiento, limpieza y transformación de datos biomédicos mediante Pandas.
    \item Comprender la estructura y el funcionamiento de bases de datos farmacológicas como DrugBank, SIDER y CIMA.
    \item Aplicar principios de ingeniería del software para desarrollar una aplicación modular, mantenible y escalable.
\end{enumerate}










